% =============================================================================
\section{Introdução}
\label{sec:introducao}

A regressão simbólica consiste em encontrar, a partir de $m$ amostras
$\langle X, Y\rangle$ com $X\in\mathbb{R}^{m\times n}$ e $Y\in\mathbb{R}^{m}$,
uma \emph{expressão} que aproxime a função desconhecida
$f:\mathbb{R}^n\to\mathbb{R}$. Ao contrário da regressão paramétrica
tradicional, não se fixa \emph{a priori} a forma do modelo: a estrutura e os
coeficientes são buscados simultaneamente. A Programação Genética
(PG)~\cite{koza1992genetic,poli2008field} é a técnica clássica para esse
problema, evoluindo populações de árvores sintáticas por seleção, cruzamento e
mutação.

O desempenho de uma PG depende de um grande número de decisões de projeto e
de parâmetros que interagem entre si. Este trabalho tem três objetivos:
(i)~implementar do zero os componentes de um arcabouço de PG para regressão
simbólica; (ii)~analisar a sensibilidade do algoritmo aos seus parâmetros
(população, gerações, seleção, operadores, elitismo e controle de tamanho);
e (iii)~avaliar as soluções de forma rigorosa, separando treino e teste,
usando 30 repetições independentes e testes estatísticos não paramétricos.

A investigação foi organizada em \textbf{duas fases}:

\begin{description}
  \item[Fase~1 --- estudo canônico.] Responde às questões do enunciado
        seguindo o guia experimental proposto: escolha de população e
        gerações, troca do método de seleção, variação de $\pc$ e $\pmut$,
        retirada do elitismo e medição do \emph{bloating}. Os melhores
        parâmetros encontrados no \code{synth1} são replicados no
        \code{synth2} e no \code{concrete1}.
  \item[Fase~2 --- estudo guiado pela \emph{fitness landscape}.] Os sintomas
        observados na Fase~1 (explosões no teste, subajuste, clones,
        crescimento das árvores, crossover destrutivo) são interpretados
        como propriedades da paisagem de aptidão, gerando seis hipóteses de
        melhoria. Cada hipótese é testada isoladamente e, depois, empilhada
        de forma auditável. Os resultados são cruzados com as perguntas da
        Fase~1 para verificar como cada intervenção altera convergência,
        diversidade, contribuição dos operadores, \emph{bloating} e
        sobreajuste.
\end{description}

As principais contribuições são: (a)~um arcabouço configurável e reprodutível
de PG com seleção por roleta, torneio e $\eps$-lexicase; (b)~um estudo
paramétrico completo com análise estatística; (c)~uma cadeia auditável
\emph{dado observado $\to$ hipótese sobre a paisagem $\to$ intervenção $\to$
veredito estatístico}; e (d)~a identificação da melhor configuração para cada
base e de uma configuração global de compromisso, com a justificativa de
\emph{por que} cada uma funciona.

O restante do artigo está estruturado conforme as diretrizes do trabalho:
a \cref{sec:implementacao} detalha a implementação completa do arcabouço,
incluindo representação em árvore, função de aptidão, operadores genéticos,
mecanismos de seleção ($\eps$-lexicase, torneio e roleta) e controle de
\emph{bloat}; a \cref{sec:experimentos} concentra a análise experimental do
impacto dos parâmetros, articulada na metodologia, no estudo canônico da
Fase~1, no aprofundamento orientado pela \emph{fitness landscape} da Fase~2, no
cruzamento de resultados e na eleição dos melhores modelos; e a
\cref{sec:conclusao} conclui o documento com as considerações finais.
